\chapter{Archie, and how these illustrations were made}
\label{app:archie}

A note that argues for archiving, referencing and reproducing research artefacts ought to say
where its own artefacts came from. Ten drawings punctuate these pages: a frontispiece and nine
vignettes, one at the head of each chapter and one over the appendices. This appendix records
what the character is, why it is an elephant, how the images were produced, and what went wrong
on the way.

\begin{callout}{Disclosure}
\small All ten illustrations are \textbf{AI-generated}, using Google Gemini's image generation,
directed and selected by the author, with light local post-processing (greyscale conversion and
white-point normalisation) in ImageMagick. No human illustrator was involved. The print copies
are in \texttt{figures/illustrations/} in the source tree of these notes, and are therefore
archived and named by the same \textsc{swhid} as the text.
\end{callout}

\section{Why a mascot}

The model is Knuth's \emph{TeXbook} and \emph{METAFONTbook}, whose whimsical pen-and-ink
vignettes, a friendly recurring lion among them, punctuate a dense technical text without
cartoonifying it. The aim here is the same: a small warm drawing at each chapter head, and one
recurring character, so that the set reads as a series rather than as a pile of clip art.

The visual style was specified by \emph{attributes}: line weight, cross-hatching, vignette
framing, whimsy. It was deliberately \emph{not} specified by naming a living illustrator, for
two reasons. Image models reproduce an aesthetic more reliably from attributes than from a name,
and naming a working artist to imitate their style is a thing one should not do.

\section{Archie, the elephant archivist}

\textbf{The character.} A gentle, rounded elephant who works as the archive's librarian: small
round spectacles, a simple archivist's apron with one breast pocket, no lettering anywhere on
the figure. Warm, unhurried, dignified. He appears in every plate, sometimes alone, once
multiplied into a network of equal peers.

\textbf{Why an elephant.} The choice is an argument in itself, and it is the book's argument.

\begin{description}[leftmargin=1.2em,itemsep=2pt,topsep=3pt]
  \item[An elephant never forgets.] The thesis of these notes is long memory: a universal
    archive that keeps the world's source code, with its full development history, for the long
    term. The proverb does the work of a paragraph.
  \item[Large, calm, dependable.] The archive is a patient, non-profit, multi-decade commitment,
    not a platform racing to a funding round. A big unhurried animal reads as durable
    custodianship rather than as move-fast development.
  \item[Expressive, and easy to pose.] Trunk, ears and spectacles allow a wide range of gentle
    business while staying dignified enough for a scholarly text. A mascot that cannot act out
    the chapter's mechanism is decoration, and decoration was not the point.
\end{description}

\textbf{Why \enquote{Archie}.} The name reads instantly as \emph{archivist} and \emph{archive},
which is the character's job and the book's subject in five letters. It also carries a quiet nod
to computing history: Archie, built at McGill University in 1990, was the first internet search
engine, the first tool made to \emph{find} what had been stored. It works unchanged in English,
French and Italian, which matters for a note that travels, and it is short enough for a caption.
Above all it names a \emph{character} rather than a concept: a recurring mascot needs a name a
reader can use.

Three other names were considered and rejected, and the reasons are instructive.

\begin{table}[htbp]
  \centering\small
  \caption[Names considered for the mascot]{Names considered for the mascot, and why each was
  set aside.}
  \label{tab:archie-names}
  \begin{tabularx}{\linewidth}{@{}>{\raggedright\arraybackslash}p{1.8cm}%
                                 >{\raggedright\arraybackslash}X%
                                 >{\raggedright\arraybackslash}X@{}}
    \toprule
    \textbf{Name} & \textbf{For} & \textbf{Against} \\
    \midrule
    \textbf{Merkle}
      & The tree he tends \emph{is} a Merkle \textsc{dag}; the name would teach the structure
        on every appearance.
      & Eponymises a living researcher, and sits oddly on a cartoon animal. \\
    \addlinespace
    \textbf{Trunk}
      & A triple pun: the elephant's trunk, the tree trunk, the version-control trunk. The
        wittiest option.
      & Names an object rather than a person, and reads awkwardly as a character name. \\
    \addlinespace
    \textbf{Mnemo}
      & From Mnemosyne; leans directly on the long-memory idea.
      & Abstract, and less pronounceable across the three languages. \\
    \bottomrule
  \end{tabularx}
\end{table}

\takeaway{A mascot earns its place when it carries the argument, not when it decorates it.}

\textbf{One standing caution.} A rounded elephant in spectacles and an apron sits visually close
to Babar. The art direction avoids names and costumes in that register, so that an incidental
family resemblance does not harden into an implied reference to a character still under
copyright. The waistcoat and necktie in the Chapter~\ref{ch:identifiers} plate are the one
departure from the apron, kept deliberately after review.

A small supporting cast appears where multiplicity is the point, and nowhere else: worker mice
with satchels in Chapter~\ref{ch:foundations}, the listers and loaders that fill the archive;
two identical round creatures at a shared drawer in the same plate, whose \emph{sameness} is the
whole joke, because identical content is stored once; and in Chapter~\ref{ch:outlook} several
small keepers, all the same size as each other, each tending an archive of their own.

\section{How the images were made}

\textbf{The tool: Google Gemini image generation.} It was chosen over the alternatives for one
specific property, \emph{character consistency across a series}. Its reference-image editing,
\enquote{here is the character, now draw him doing this}, is what keeps Archie recognisably the
same elephant across ten drawings. The trade-off is weaker literal adherence to dense
compositional prose, which is exactly where the failures in the next section came from.

\textbf{The consistency chain.} The images were not generated independently. They were chained:

\begin{enumerate}[leftmargin=1.4em,itemsep=2pt,topsep=3pt]
  \item Generate one vignette first with no reference, and regenerate it until the design is
    right. It establishes the canonical Archie, and becomes the reference sheet.
  \item For every later image, \emph{attach that reference}, instruct the model to keep the
    character identical, and describe only the new scene. Attaching the image is what carries
    the character; words alone do not.
  \item On drift, re-attach the best prior image: \enquote{match this elephant exactly, same
    face, spectacles, apron; change only the scene.}
  \item State the aspect ratio in words (\enquote{landscape 3:2 composition}), which the model
    honours.
  \item Regenerate on stray lettering rather than keeping it. Garbled characters read as an
    error in a scholarly book.
\end{enumerate}

\noindent A fixed style anchor was prepended to every scene prompt, and never varied:

\begin{quote}\small
\textbf{Style:} black-and-white \textbf{pen-and-ink line drawing}, confident hand-drawn linework
of varying weight, light \textbf{cross-hatching} for shading, generous white space, no fill tones
or grey washes beyond hatching. A small, warm \textbf{vignette}, a single self-contained little
scene that fades softly at its edges into the white of the page, \emph{not} a full bordered
illustration and \emph{not} a busy full scene. The mood is whimsical, gentle and good-humoured,
in the tradition of a mid-twentieth-century technical-book chapter-head drawing.
\textbf{No colour. No photorealism, no digital gloss, no gradients, no 3D rendering. No text, no
lettering, no numbers, no logos anywhere in the image.} Recurring character: a gentle, rounded
\textbf{elephant archivist} with small round spectacles and a simple archivist's apron, drawn the
same way each time.
\end{quote}

\section{What each plate says}

Each vignette states its chapter's mechanism, not its mood.

\begin{description}[leftmargin=1.2em,itemsep=2pt,topsep=3pt,font=\normalfont\itshape]
  \item[Frontispiece.] A tree rooted in an open book, its leaves made of files, with Archie
    steadying a branch: the science grows out of its source.
  \item[Front matter.] Three pedestals bearing a bound book, a specimen jar and a tree-shaped
    bundle of files, the trunk gesturing to the third as to an old friend left out of things.
    Publications, data and software are three research outputs of equal standing.
  \item[Chapter~\ref{ch:why}.] The journal is shelved on solid ground while the code tumbles off
    a crumbling ledge into loose leaves, just past the reach of the outstretched trunk. One
    output is preserved by an apparatus built for it; the other is not.
  \item[Chapter~\ref{ch:beyond-fair}.] A sealed jar with a pinned specimen, held up against a
    potted sapling that sprouts, reaches tendrils to its neighbours and overgrows its pot. Data
    holds still; software does not.
  \item[Chapter~\ref{ch:foundations}.] A tree built of shelves and drawers, files at the tips,
    gathered into branches, gathered into one trunk, with one branch growing \emph{into} the
    branch above it, because the archive is a directed acyclic graph and a subtree can have
    several parents.
  \item[Chapter~\ref{ch:identifiers}.] The same bundle named twice: a paper luggage tag on a
    string running off to a distant registry desk, eyed with scepticism, against a fingerprint
    pressed into the object's own surface and read on the spot through a lens.
  \item[Chapter~\ref{ch:ardc}.] The right hand writes in the open citation catalogue; the left
    holds a lens flat over a completely blank book, and the fingerprint appears only inside the
    glass. The identifier is computed from the content and then transcribed, never stamped on.
  \item[Chapter~\ref{ch:repro}.] A fair examiner at a desk, matching a submitted bundle's mark
    against the claimed mark on a slip, a blank ribboned rosette ready to award, a folding screen
    hiding who sent it. Verification needs no authority, and survives double-blind review.
  \item[Chapter~\ref{ch:outlook}.] A globe of plain line-work continents, several equal keepers
    tending their own patch, their sites linked by faint dotted arcs, one of them offering a
    sapling toward the reader. A mirror network has no single point of failure and no hero.
  \item[Appendices.] Archie in a cook's apron, spoon behind one ear, pulling a card from a recipe
    box: what follows is a cookbook, recipes to follow rather than arguments to weigh. It is the
    one image that deliberately contains lettering, because the words are the metaphor.
\end{description}

\section{What went wrong, and what fixed it}

Three failures are worth recording, because they generalise beyond drawings.

\textbf{Name the delta; never re-describe the scene.} A prompt that restates a scene is read as a
request to regenerate it, not to edit it. Every style clause repeated for safety is
re-interpreted from scratch, and that, rather than poor taste, produced every drift observed
here: sepia backgrounds, an engraving style, four elephants where one was wanted. The reliable
form is a preservation clause and one sentence of change, naming the region by landmark:
\enquote{keep this exact image unchanged; change only \textellipsis}.

\textbf{A drawing can assert the opposite of the text.} An early version of the
Chapter~\ref{ch:ardc} plate had Archie \emph{stamping} an identifier onto a citation slip. That
is the extrinsic gesture Chapter~\ref{ch:identifiers} argues against, drawn at the head of the
chapter that depends on the distinction, and it was discarded for that reason rather than
retouched. In the same plate the model kept printing the fingerprint on the book cover as well as
in the lens, which quietly restores the very claim the picture exists to refute. An illustration
is part of the argument and has to be reviewed as such.

\textbf{Judge the artefact, not its description.} The Chapter~\ref{ch:identifiers} plate carried
a small hanging tag reading \texttt{LUGGAGE LABL}: legible, unintended, misspelled, and forbidden
by every prompt in the set. It survived an author-side review conducted from the prompts, the
thumbnails and the file names, and was caught only when a reviewer opened each image at full
resolution with no expectation of what it should contain. The lettering was then removed by a
local edit of a few thousand pixels rather than by regeneration, which would have risked
character drift for a defect confined to one tag.

\takeaway{A name, a prompt or a filename describes an artefact. Only the artefact is evidence.}

That last failure is this book's own thesis turned back on its own production. An
identifier, a description, a record of intent: none of them is the object. The reason to name
source code by a hash of its content is precisely that the hash is computed from what is there,
and not from what someone meant to put there.
